\Opensolutionfile{ans}[ans/ansTN-BTX-429]

\noindent \textbf{A. Phần trắc nghiệm (7,0 điểm)}

\noindent \textbf{PHẦN I. Câu hỏi có nhiều phương án lựa chọn. (3,0đ) Thí sinh trả lời từ câu 1 đến câu 12. Mỗi câu hỏi thí sinh chỉ chọn một phương án}
\vspace{0.2cm}

\begin{ex}	\macau{ID: VL12-BTX429-C01}Cồn y tế chuyển từ thể lỏng sang thể khí rất nhanh ở điều kiện thông thường. Khi xoa cồn vào da, ta cảm thấy lạnh ở vùng da đó vì cồn
	\choice
	{khi bay hơi kéo theo lượng nước chỗ da đó ra khỏi cơ thể}
	{khi bay hơi tỏa nhiệt lượng vào chỗ da đó}
	{khi bay hơi tạo ra dòng nước mát tại chỗ da đó}
	{\True thu nhiệt lượng từ cơ thể qua chỗ da đó để bay hơi}
	\loigiai{
		Quá trình chuyển pha từ thể lỏng sang thể khí (bay hơi) là quá trình thu nhiệt. Khi xoa cồn lên da, cồn thu nhiệt lượng trực tiếp từ cơ thể qua vùng da đó để cung cấp cho quá trình bay hơi, làm cho nhiệt độ vùng da đó giảm xuống nên ta cảm thấy lạnh}
\end{ex}

\begin{ex}	\macau{ID: VL12-BTX429-C02}Tính chất nào sau đây \textbf{không} phải là tính chất của chất ở thể khí?
	\choice
	{Có thể nén được dễ dàng}
	{Có lực tương tác phân tử nhỏ hơn lực tương tác phân tử ở thể rắn và thể lỏng}
	{\True Có hình dạng và thể tích riêng xác định}
	{Có các phân tử chuyển động hoàn toàn hỗn độn}
	\loigiai{
		Chất ở thể khí có lực tương tác giữa các phân tử rất yếu, các phân tử chuyển động hỗn loạn không ngừng chiếm toàn bộ thể tích của bình chứa nên chất khí không có hình dạng và thể tích riêng xác định. Nhận định C sai}
\end{ex}

\begin{ex}	\macau{ID: VL12-BTX429-C03}Công thức nào sau đây mô tả đúng định luật I (Nguyên lí I) của nhiệt động lực học?
	\choice
	{$\Delta U=A-Q$}
	{\True $\Delta U=A+Q$}
	{$\Delta U=Q-A$}
	{$A=\Delta U-Q$}
	\loigiai{
		Biểu thức chuẩn xác của Nguyên lí I nhiệt động lực học mô tả sự biến thiên nội năng của hệ bằng tổng công và nhiệt lượng mà hệ trao đổi với bên ngoài là: $\Delta U = Q + A$}
\end{ex}

\begin{ex}	\macau{ID: VL12-BTX429-C04}Cách xác định nhiệt độ trong thang nhiệt độ Celsius là
	\choice
	{Lấy nhiệt độ của nước khi đóng băng là ($10^\circ\text{C}$) và nhiệt độ sôi của nước ($100^\circ\text{C}$) làm chuẩn}
	{Lấy nhiệt độ của nước khi đóng băng là ($100^\circ\text{C}$) và nhiệt độ sôi của nước ($10^\circ\text{C}$) làm chuẩn}
	{\True Lấy nhiệt độ của nước khi đóng băng là ($0^\circ\text{C}$) và nhiệt độ sôi của nước làm chuẩn ($100^\circ\text{C}$)}
	{Lấy nhiệt độ của nước khi đóng băng là ($10^\circ\text{C}$) và nhiệt độ sôi của nước ($0^\circ\text{C}$) làm chuẩn}
	\loigiai{
		Nhiệt giai bách phân Celsius chọn hai mốc cố định: mốc nhiệt độ đóng băng của nước tinh khiết là $0^\circ\text{C}$ và mốc nhiệt độ hơi nước đang sôi bão hòa dưới áp suất tiêu chuẩn là $100^\circ\text{C}$}
\end{ex}

\begin{ex}	\macau{ID: VL12-BTX429-C05}Quy ước dấu nào sau đây phù hợp với Nguyên lí I của nhiệt động lực học?
	\choice
	{\True Vật nhận công: $A>0$; vật nhận nhiệt lượng: $Q>0$}
	{Vật nhận công: $A<0$; vật nhận nhiệt lượng: $Q<0$}
	{Vật thực hiện công: $A>0$; vật truyền nhiệt lượng: $Q>0$}
	{Vật thực hiện công: $A>0$; vật truyền nhiệt lượng: $Q<0$}
	\loigiai{
		Quy ước dấu tiêu chuẩn của Nguyên lí I nhiệt động lực học quy định: Hệ nhận nhiệt lượng từ bên ngoài $\Rightarrow Q > 0$; Hệ nhận công từ ngoại lực tác dụng $\Rightarrow A > 0$}
\end{ex}

\begin{ex}	\macau{ID: VL12-BTX429-C06}Nhiệt dung riêng của sắt là $440\text{ J/(kg.K)}$. Điều này có nghĩa là
	\choice
	{nếu lấy đi nhiệt lượng $440\text{ J}$ thì nhiệt độ của $1\text{ kg}$ sắt sẽ tăng thêm $1^\circ\text{C}$}
	{để làm cho $1\text{ kg}$ sắt tăng nhiệt độ từ $0^\circ\text{C}$ đến $100^\circ\text{C}$ cần $440\text{ J}$}
	{để làm nóng chảy $1\text{ kg}$ sắt cần $440\text{ J}$}
	{\True nếu lấy đi nhiệt lượng $440\text{ J}$ thì nhiệt độ của $1\text{ kg}$ sắt sẽ giảm đi $1^\circ\text{C}$}
	\loigiai{
		Nhiệt dung riêng $c = 440\text{ J/(kg.K)}$ cho biết để làm cho $1\text{ kg}$ sắt thay đổi nhiệt độ một lượng $1\text{ K}$ (hoặc $1^\circ\text{C}$) thì cần cung cấp (nếu tăng nhiệt độ) hoặc lấy đi (nhếu giảm nhiệt độ) một nhiệt lượng bằng đúng $440\text{ J}$. Phương án D diễn đạt hoàn toàn chính xác}
\end{ex}

\begin{ex}	\macau{ID: VL12-BTX429-C07}Hãy chỉ ra phương án \textbf{sai} trong các câu sau: Cùng một khối lượng của một chất nhưng khi ở các thể khác nhau thì sẽ khác nhau về
	\choice
	{trật tự sắp xếp của các nguyên tử}
	{\True kích thước hình học của bản thân các nguyên tử}
	{thể tích vĩ mô của khối chất}
	{khối lượng riêng của chất}
	\loigiai{
		Khi một chất chuyển pha cấu trúc ở các thể khác nhau (rắn, lỏng, khí), sự khác biệt nằm ở khoảng cách và trật tự sắp xếp phân tử, kéo theo thể tích và khối lượng riêng thay đổi. Bản thân kích thước lõi của từng nguyên tử, phân tử cấu cấu tử là hoàn toàn cố định không đổi. Phương án B là phương án sai}
\end{ex}

\begin{ex}	\macau{ID: VL12-BTX429-C08}Nội năng của một vật được định nghĩa là
	\choice
	{\True tổng động năng chuyển động nhiệt và thế năng tương tác của các phân tử cấu tạo nên vật}
	{tổng nhiệt lượng và cơ năng mà vật nhận được trong quá trình truyền nhiệt và thực hiện công}
	{nhiệt lượng vật nhận được riêng trong quá trình truyền nhiệt}
	{tổng động năng và thế năng cơ học vĩ mô của vật}
	\loigiai{
		Nội năng ($U$) là một hàm trạng thái năng lượng bên trong của hệ vật chất, bằng tổng động năng chuyển động nhiệt hỗn loạn của các hạt cộng với thế năng tương tác tương hỗ giữa các phân tử cấu tạo nên vật đó}
\end{ex}

\begin{ex}	\macau{ID: VL12-BTX429-C09}Trong hệ thống SI, đơn vị tiêu chuẩn dùng để đo hằng số nhiệt nóng chảy riêng $\lambda$ của khối chất rắn là
	\choice
	{$\text{J/(kg.K)}$}
	{\True $\text{J/kg}$}
	{$\text{J/s}$}
	{$\text{kg/J}$}
	\loigiai{
		Nhiệt nóng chảy riêng được xác định từ công thức chuyển pha: $\lambda = \frac{Q}{m}$. Đơn vị đo tương ứng tương ứng là Jun trên kilôgam ($\text{J/kg}$)}
\end{ex}

\begin{ex}	\macau{ID: VL12-BTX429-C10}Các dụng cụ đo nhiệt kế chất lỏng thông dụng được chế tạo vận hành dựa trên nguyên tắc vật lí nào?
	\choice
	{\True Sự nở vì nhiệt tự nhiên của các chất lỏng}
	{Sự nở ra của chất lỏng khi nhiệt độ bị sụt giảm}
	{Sự co rút lại của chất lỏng khi nhiệt độ dâng tăng}
	{Sự nở của chất lỏng hoàn toàn không phụ thuộc vào nhiệt độ}
	\loigiai{
		Nhiệt kế chất lỏng (như nhiệt kế thủy ngân, nhiệt kế rượu) vận hành đo đạc dựa trên thuộc tính dãn nở dài và dãn nở thể tích tự nhiên vì nhiệt của chất lỏng mao dẫn chứa ở bầu đo}
\end{ex}

\begin{ex}	\macau{ID: VL12-BTX429-C11}Trị số điểm thềm nhiệt độ "Độ không tuyệt đối" ứng với giá trị nào sau đây trên các thang đo?
	\choice
	{\True $0\text{ K}$}
	{$273^\circ\text{C}$}
	{$73\text{ K}$}
	{$0^\circ\text{C}$}
	\loigiai{
		Độ không tuyệt đối là mốc giới hạn nhiệt độ thấp nhất trên lý thuyết, ứng với mốc vạch $0\text{ K}$ trên thang đo nhiệt tuyệt đối Kelvin (tương đương $-273,15^\circ\text{C}$ Thang Celsius)}
\end{ex}

\begin{ex}	\macau{ID: VL12-BTX429-C12}Nhiệt hoá hơi riêng của một chất là lượng nhiệt lượng cần thiết cung cấp để cho $1\text{ kg}$ chất đó
	\choice
	{hoá hơi hoàn toàn}
	{bay hơi hết ở mọi điều kiện}
	{\True hoá hơi hoàn toàn ở đúng vạch điểm nhiệt độ sôi}
	{bắt đầu quá trình hoá hơi}
	\loigiai{
		Theo định nghĩa, nhiệt hóa hơi riêng $L$ là nhiệt lượng cần cung cấp để làm cho một đơn vị khối lượng ($1\text{ kg}$) chất lỏng chuyển pha hóa hơi hoàn toàn ổn định ở đúng nhiệt độ sôi xác định dưới áp suất cho trước}
\end{ex}

\Closesolutionfile{ans}

%%%%%%%%%%%%%%%%%%%%%%%%%%%%%%%%%%%%%%%%%%%%%%%%%%%%%%%%%%%%%%%%%%%%%%%%%
% PHẦN II: CÂU HỎI TRẮC NGHIỆM ĐÚNG SAI (2 CÂU)
%%%%%%%%%%%%%%%%%%%%%%%%%%%%%%%%%%%%%%%%%%%%%%%%%%%%%%%%%%%%%%%%%%%%%%%%%
\noindent \textbf{PHẦN II. Câu hỏi trắc nghiệm đúng sai. (2,0đ) Thí sinh trả lời từ câu 1 đến câu 2. Trong mỗi ý a), b), c), d) ở mỗi câu, thí sinh chọn đúng hoặc sai}
\Opensolutionfile{ans}[ans/ansTF-BTX-429]

\begin{ex}	\macau{ID: VL12-BTX429-C13}Người ta thực hiện một công cơ học có độ lớn bằng $100\text{ J}$ để đè nén lượng chất khí chứa trong một bình xi-lanh kín, đồng thời khối chất khí này truyền tỏa ra môi trường xung quanh một nhiệt lượng có độ lớn bằng $20\text{ J}$.
	\choiceTF[t]
	{Nội năng của khối khí chứa trong xi-lanh chỉ có thể bị thay đổi bằng phương thức duy nhất là thực hiện công}
	{\True Do khối khí thực hiện truyền nhiệt lượng ra môi trường ngoài nên giá trị đại số quy ước dấu mang trị số $Q = -20\text{ J}$}
	{Trong quá trình nén ép cơ học trên, khối chất khí đã nhận được một lượng công đại số mang trị số $A = +100\text{ J}$}
	{\True Tổng hợp toàn bộ quá trình biến đổi trên, giá trị nội năng của khối khí lỏng đã được gia tăng thêm một lượng bằng $80\text{ J}$}
	\loigiai{
		\begin{itemchoice}
			\itemch Nội năng của một hệ nhiệt động có thể bị biến đổi bằng hai phương thức vĩ mô độc lập: thực hiện công cơ học hoặc truyền nhiệt năng thuần túy.
			\itemch Theo quy ước dấu Nguyên lí I, hệ tỏa nhiệt lượng (truyền nhiệt) ra ngoài môi trường nhận dấu âm: $Q = -20\text{ J}$.
			\itemch Khí chịu nén từ lực ngoại cảnh dập xuống tức là khí nhận công, do đó hằng số công nhận dấu dương đại số là $A = +100\text{ J}$. Giá trị $120\text{ J}$ ở mệnh đề là sai.
			\itemch Áp dụng hệ thức định luật Nguyên lí I nhiệt động lực học tính độ biến thiên nội năng:
			$$\Delta U = Q + A = -20 + 100 = +80\text{ J}$$
			Vì $\Delta U > 0$ nên nội năng của khối khí tăng một lượng bằng $80\text{ J}$.
		\end{itemchoice}
	}
\end{ex}

\begin{ex}	\macau{ID: VL12-BTX429-C14}Thả một cục nước đá có khối lượng $50\text{ g}$ đang ở nhiệt độ nóng chảy $0^\circ\text{C}$ vào một chiếc cốc chứa khối lượng nước lỏng bằng $200\text{ g}$ đang ở vạch nhiệt độ $25^\circ\text{C}$. Giả định hệ cô lập hoàn toàn bỏ qua mọi sự trao đổi nhiệt với vỏ cốc và môi trường không khí ngoài. Biết ẩn nhiệt nóng chảy riêng của nước đá cho bằng $\lambda = 3,4 \cdot 10^5\text{ J/kg}$, và hằng số nhiệt dung riêng của nước lỏng bằng $c = 4180\text{ J/(kg.K)}$.
	\choiceTF[t]
	{Ngay sau khi thả cục nước đá vào, dòng năng lượng nhiệt tự phát đã thực hiện dịch chuyển truyền thẳng từ cục đá sang khối nước lỏng}
	{\True Tổng lượng nhiệt lượng do khối nước trong cốc tỏa ra khi hệ thiết lập trạng thái cân bằng nhiệt bằng đúng lượng nhiệt cục nước đá thu vào để tan và tăng nhiệt}
	{\True Nhiệt lượng cần thiết cung cấp để làm tan chảy hoàn toàn cục nước đá trên ở mức nhiệt độ nóng chảy bằng $17000\text{ J}$}
	{\True Mức vạch nhiệt độ cuối cùng khi hệ đạt trạng thái ổn định cân bằng nhiệt bấy giờ đạt trị số xấp xỉ bằng $3,7^\circ\text{C}$}
	\loigiai{
		\begin{itemchoice}
			\itemch Chiều tự phát của dòng truyền nhiệt luôn dịch chuyển từ vật có vạch nhiệt độ cao hơn (nước lỏng ấm $25^\circ\text{C}$) sang vật có nhiệt độ thấp hơn (nước đá lạnh $0^\circ\text{C}$).
			\itemch Do hệ kín cô lập không thất thoát năng lượng ra ngoài, phương trình bảo toàn năng lượng chỉ ra tổng nhiệt tỏa ra bằng tổng nhiệt thu vào: $Q_{\text{tỏa}} = Q_{\text{thu}}$.
			\itemch Quy đổi khối lượng đá: $m_1 = 50\text{ g} = 0,05\text{ kg}$. Lượng nhiệt phá vỡ liên kết tan chảy đá hoàn toàn: 
			$$Q_{\text{nc}} = \lambda \cdot m_1 = 3,4 \cdot 10^5 \cdot 0,05 = 17000\text{ J}$$
			\itemch Nhiệt lượng tối đa khối nước ấm tỏa ra khi hạ hẳn nhiệt độ về thềm đá tan $0^\circ\text{C}$ là: $Q_{\text{tỏamax}} = m_2 \cdot c \cdot (t_2 - 0) = 0,2 \cdot 4180 \cdot 25 = 20900\text{ J}$. \\
			Vì $Q_{\text{tỏamax}} = 20900\text{ J} > Q_{\text{nc}} = 17000\text{ J}$ nên nước đá bị tan chảy hoàn toàn hóa lỏng và lượng dư tiếp tục làm ấm hệ lên đến nhiệt độ cân bằng chung $t_{\text{cb}} > 0$. Lập phương trình cân bằng nhiệt toàn phần:
			$$Q_{\text{tỏa}} = Q_{\text{thu}} \Rightarrow m_2 \cdot c \cdot (t_2 - t_{\text{cb}}) = Q_{\text{nc}} + m_1 \cdot c \cdot (t_{\text{cb}} - 0)$$
			$$0,2 \cdot 4180 \cdot (25 - t_{\text{cb}}) = 17000 + 0,05 \cdot 4180 \cdot t_{\text{cb}}$$
			$$20900 - 836t_{\text{cb}} = 17000 + 209t_{\text{cb}} \Rightarrow 3900 = 1045t_{\text{cb}} \Rightarrow t_{\text{cb}} \approx 3,73^\circ\text{C} \approx 3,7^\circ\text{C}$$
		\end{itemchoice}
	}
\end{ex}

\Closesolutionfile{ans}

%%%%%%%%%%%%%%%%%%%%%%%%%%%%%%%%%%%%%%%%%%%%%%%%%%%%%%%%%%%%%%%%%%%%%%%%%
% PHẦN III: CÂU HỎI TRẮC NGHIỆM TRẢ LỜI NGẮN
%%%%%%%%%%%%%%%%%%%%%%%%%%%%%%%%%%%%%%%%%%%%%%%%%%%%%%%%%%%%%%%%%%%%%%%%%
\noindent \textbf{PHẦN III. Câu hỏi trắc nghiệm trả lời ngắn (2,0đ)}

\noindent \textit{(Thí sinh trả lời từ câu 1 đến câu 4)}
\vspace{0.2cm}

\begin{ex}	\macau{ID: VL12-BTX429-C15}Nội năng của một khối khí lí tưởng gia tăng thêm một lượng bằng $10\text{ J}$ khi người ta tiến hành truyền cấp cho khối khí này một lượng nhiệt lượng bằng $40\text{ J}$. Hỏi trong suốt tiến trình biến đổi nhiệt động đó, khối khí đã thực hiện sinh ra một công cơ học có độ lớn bằng bao nhiêu Jun?

	\par\shortans[oly]{30}
	\loigiai{
		\begin{itemize}
			\item Trích xuất số liệu đại số theo quy ước dấu: $\Delta U = +10\text{ J}$ (nội năng tăng); $Q = +40\text{ J}$ (hệ nhận nhiệt lượng).
			\item Áp dụng công thức hệ thức Nguyên lý I nhiệt động lực học để tìm giá trị công $A$:
			$$\Delta U = Q + A \Rightarrow +10 = 40 + A \Rightarrow A = 10 - 40 = -30\text{ J}$$
			\item Giá trị công đại số $A = -30\text{ J}$ cho biết khối khí sinh công (thực hiện công) ra bên ngoài với độ lớn công cơ học tương ứng bằng $|A| = 30\text{ J}$.
		\end{itemize}
	}
\end{ex}

\begin{ex}	\macau{ID: VL12-BTX429-C16}Một vật thể được làm lạnh sâu hạ nhiệt độ từ vạch mức $25^\circ\text{C}$ xuống đến vạch mức hạ lưu bằng $3^\circ\text{C}$. Hỏi vạch nhiệt độ của vật thể đó dóng theo thang đo tuyệt đối Kelvin đã sụt giảm đi một lượng bằng bao nhiêu Kelvin (K)?

	\par\shortans[oly]{22}
	\loigiai{
		\begin{itemize}
			\item Độ chênh lệch biến thiên nhiệt độ tính theo thang bách phân Celsius: $\Delta t = 25 - 3 = 22^\circ\text{C}$.
			\item Vì độ rộng của một độ chia vạch trong thang đo bách phân Celsius hoàn toàn trùng khít bằng nhau so với độ rộng của một độ chia trong thang tuyệt đối Kelvin, nên lượng biến thiên nhiệt độ ở hai thang đo luôn trọn vẹn bằng nhau: $\Delta T = \Delta t = 22\text{ K}$.
		\end{itemize}
	}
\end{ex}

\begin{ex}	\macau{ID: VL12-BTX429-C17}Tính tổng nhiệt lượng tối thiểu cần cung cấp để chuyển một khối lượng $400\text{ g}$ nước đá từ mức nhiệt độ ban đầu $-5^\circ\text{C}$ chuyển pha biến đổi hóa hơi hoàn toàn ở điểm sôi bão hòa $100^\circ\text{C}$. Quy đổi kết quả đo theo đơn vị đo triệu Jun (MJ) và làm tròn lấy đến hai chữ số sau dấu phẩy thập phân? Cho biết các thông số hằng số nhiệt kĩ thuật: nhiệt dung riêng của đá bằng $2100\text{ J/(kg.K)}$, của nước lỏng bằng $4200\text{ J/(kg.K)}$; ẩn nhiệt nóng chảy riêng của đá bằng $3,34 \cdot 10^5\text{ J/kg}$ và hằng số ẩn nhiệt hóa hơi riêng của nước ở điểm sôi bằng $2,26 \cdot 10^6\text{ J/kg}$.

	\par\shortans[oly]{1,21}
	\loigiai{
		\begin{itemize}
			\item Quy đổi khối lượng mẫu chất sang kg chuẩn: $m = 400\text{ g} = 0,4\text{ kg}$. Tiến trình thu nhiệt gồm 4 giai đoạn nối tiếp:
			\begin{enumerate}
				\itemch Giai đoạn 1: Thu nhiệt để nâng đá từ $-5^\circ\text{C} \rightarrow 0^\circ\text{C}$: $Q_1 = m \cdot c_{\text{đá}} \cdot \Delta t_1 = 0,4 \cdot 2100 \cdot [0 - (-5)] = 4200\text{ J}$.
				\itemch Giai đoạn 2: Thu nhiệt nóng chảy tan đá hoàn toàn ở $0^\circ\text{C}$: $Q_2 = m \cdot \lambda = 0,4 \cdot 3,34 \cdot 10^5 = 133600\text{ J}$.
				\itemch Giai đoạn 3: Thu nhiệt nâng nước lỏng từ $0^\circ\text{C} \rightarrow 100^\circ\text{C}$: $Q_3 = m \cdot c_{\text{nước}} \cdot \Delta t_3 = 0,4 \cdot 4200 \cdot (100 - 0) = 168000\text{ J}$.
				\itemch Giai đoạn 4: Thu ẩn nhiệt để hóa hơi nước bão hòa hoàn toàn ở $100^\circ\text{C}$: $Q_4 = m \cdot L = 0,4 \cdot 2,26 \cdot 10^6 = 904000\text{ J}$.
			\end{enumerate}
			\item Tổng nhiệt lượng toàn phần tiêu hao trên cả hành trình:
			$$Q_{\text{tổng}} = Q_1 + Q_2 + Q_3 + Q_4 = 4200 + 133600 + 168000 + 904000 = 1209800\text{ J} = 1,2098\text{ MJ}$$
			\item Làm tròn lấy hai chữ số thập phân theo yêu cầu thu được kết quả chỉ số: $1,21\text{ MJ}$.
		\end{itemize}
	}
\end{ex}

\begin{ex}	\macau{ID: VL12-BTX429-C18}Một người quản lí muốn chuẩn bị nước pha bồn tắm hỗn hợp đạt mức vạch nhiệt độ ổn định bằng $35^\circ\text{C}$ bằng cách kết hợp đổ dòng nước sôi ($100^\circ\text{C}$) hòa chung vào một chiếc thùng đang chứa sẵn $6,2\text{ lít}$ nước lạnh đang ở vạch nhiệt độ $25^\circ\text{C}$. Giả định hệ lý tưởng không có hao phí nhiệt lượng truyền ra ngoài vỏ bình. Hãy tính toán xác định lượng dung tích nước sôi tối thiểu cần đun rót thêm vào bằng bao nhiêu lít?

	\par\shortans[oly]{0,95}
	\loigiai{
		\begin{itemize}
			\item Vì khối lượng riêng của nước lỏng ổn định ($\rho \approx 1\text{ kg/lít}$), lượng dung tích $6,2\text{ lít}$ nước lạnh tương ứng khối lượng $m_2 = 6,2\text{ kg}$. Gọi $m_1$ là khối lượng nước sôi cần rót.
			\item Lập phương trình cân bằng nhiệt lượng thu tỏa của quá trình trộn chất lưu:
			$$Q_{\text{tỏa}} = Q_{\text{thu}} \Rightarrow m_1 \cdot c \cdot (t_{\text{sôi}} - t_{\text{cb}}) = m_2 \cdot c \cdot (t_{\text{cb}} - t_{\text{lạnh}})$$
			$$m_1 \cdot (100 - 35) = 6,2 \cdot (35 - 25) \Rightarrow 65 \cdot m_1 = 6,2 \cdot 10 = 62$$
			$$m_1 = \frac{62}{65} \approx 0,9538\text{ kg}$$
			\item Quy đổi giá trị khối lượng sang vạch chỉ số dung tích lít tương ứng: $V_1 \approx 0,95\text{ lít}$.
		\end{itemize}
	}
\end{ex}

\Closesolutionfile{ans}

%%%%%%%%%%%%%%%%%%%%%%%%%%%%%%%%%%%%%%%%%%%%%%%%%%%%%%%%%%%%%%%%%%%%%%%%%
% B. PHẦN TỰ LUẬN TỰ LUẬN BÀI TẬP (3 ĐIỂM)
%%%%%%%%%%%%%%%%%%%%%%%%%%%%%%%%%%%%%%%%%%%%%%%%%%%%%%%%%%%%%%%%%%%%%%%%%
\noindent \textbf{B. Phần tự luận (3,0đ)}
\vspace{0.2cm}

\begin{ex}	\macau{ID: VL12-BTX429-C19}\textbf{(1,0 điểm)} Người ta tiến hành cấp một nhiệt lượng có độ lớn bằng $3,5\text{ J}$ truyền vào cho một khối khí lí tưởng chứa bên trong một bình xi-lanh đặt nằm ngang. Khối chất khí hấp thụ năng lượng dãn nở đẩy tấm pít-tông dịch chuyển chuyển động tịnh tiến đi được một khoảng hành trình dài đúng $6\text{ cm}$. Biết lực ma sát trượt cơ học xuất hiện ở bề mặt tiếp xúc giữa pít-tông và thành vách xi-lanh có độ lớn bằng $25\text{ N}$ và coi pít-tông chuyển động thẳng đều. Hãy tính toán xác định:
	\begin{enumerate}
		\item[a.] Độ lớn phần công cơ học do khối chất khí đã thực hiện sinh ra.
		\item[b.] Trị số độ biến thiên nội năng $\Delta U$ đại số của khối khí khí.
	\end{enumerate}
	\loigiai{
		\begin{enumerate}
			\item[a.] Quy đổi khoảng hành trình dịch chuyển sang mét chuẩn hệ SI: $s = 6\text{ cm} = 0,06\text{ m}$. \\
			Vì tấm pít-tông thực hiện chuyển động thẳng đều, lực đẩy cơ học do áp suất khối khí sinh ra tác dụng lên pít-tông cân bằng động học hoàn toàn về mặt độ lớn đối với lực ma sát trượt cản phá: $F_{\text{đẩy}} = F_{\text{ms}} = 25\text{ N}$. \\
			Độ lớn công cơ học do khối chất khí thực hiện sinh ra đẩy pít-tông dời chỗ bằng:
			$$A' = F_{\text{đẩy}} \cdot s = 25 \cdot 0,06 = 1,5\text{ J}$$
			Hệ thực hiện công (sinh công) nên hằng số công đại số nhận giá trị âm tương ứng: $A = -1,5\text{ J}$.
			\item[b.] Trích xuất thông số số đo nhiệt lượng hệ nhận vào từ nguồn ngoài: $Q = +3,5\text{ J}$. \\
			Áp dụng hệ thức toán học định luật Nguyên lí I nhiệt động lực học để tính biến thiên nội năng:
			$$\Delta U = Q + A = 3,5 + (-1,5) = +2,0\text{ J}$$
			Vì $\Delta U > 0$, nội năng của khối khí tăng một lượng bằng đúng $2,0\text{ J}$.
		\end{enumerate}
	}
\end{ex}

\begin{ex}	\macau{ID: VL12-BTX429-C20}\textbf{(1,0 điểm)} Một lò thổi đúc luyện lò cao của một nhà máy luyện kim công nghiệp có công suất lớn, mỗi lần vận hành thực hiện nung đúc chảy định mức được khối lượng bằng $30\text{ tấn}$ thép thành phẩm. Cho biết thông số ẩn nhiệt nóng chảy riêng của vật liệu thép hợp kim lấy bằng $\lambda = 2,77 \cdot 10^5\text{ J/kg}$.
	\begin{enumerate}
		\item[a.] Tính tổng nhiệt lượng định mức cần thiết đun cấp để làm nóng chảy hóa lỏng hoàn toàn khối lượng thép trong mỗi lần luyện của nhà máy tại đúng điểm thềm nhiệt độ nóng chảy của nó.
		\item[b.] Giả sử hệ thống lò cao sử dụng nguồn năng lượng đốt bằng khí hóa lỏng để nấu chảy thép. Biết kĩ thuật chỉ rõ khi đốt cháy hoàn toàn $1\text{ kg}$ khí đốt chuyên dụng này thì lượng nhiệt lượng tỏa ra ngoài môi trường đạt mức hằng số bằng $q = 45 \cdot 10^5\text{ J}$. Xác định khối lượng khí đốt tiêu hao tối thiểu cần phải sử dụng cho mỗi lượt luyện thép của nhà máy, coi hiệu suất chuyển hóa nhiệt lượng lò đạt mức lý tưởng lý tưởng.
	\end{enumerate}
	\loigiai{
		\begin{enumerate}
			\item[a.] Quy đổi khối lượng vật liệu thép sang hệ tấn sang đơn vị kilôgam tiêu chuẩn SI: 
			$$m = 30\text{ tấn} = 30000\text{ kg} = 3 \cdot 10^4\text{ kg}$$
			Nhiệt lượng chuyển pha nóng chảy cần cung cấp để làm nóng chảy hoàn toàn khối lượng thép ở thềm điểm nhiệt độ nóng chảy xác định bằng biểu thức:
			$$Q = \lambda \cdot m = 2,77 \cdot 10^5\text{ J/kg} \cdot 3 \cdot 10^4\text{ kg} = 8,31 \cdot 10^9\text{ J}$$
			\item[b.] Áp dụng phương trình định luật bảo toàn năng lượng hóa năng tỏa ra khi đốt cháy nhiên liệu, lượng nhiệt lượng tỏa ra khi đốt cháy khối lượng $m_{\text{khí}}$ khí hóa lỏng bằng: $Q_{\text{tỏa}} = q \cdot m_{\text{khí}}$. \\
			Do hệ thống vận hành lý tưởng lý tưởng, toàn bộ nhiệt lượng tỏa ra cấp trọn vẹn cho khối thép chuyển pha: $Q_{\text{tỏa}} = Q$. \\
			Khối lượng khí đốt hóa lỏng tiêu hao cần tiêu thụ bằng:
			$$m_{\text{khí}} = \frac{Q}{q} = \frac{8,31 \cdot 10^9\text{ J}}{45 \cdot 10^5\text{ J/kg}} \approx 1846,67\text{ kg}$$
		\end{enumerate}
	}
\end{ex}

\begin{ex}	\macau{ID: VL12-BTX429-C21}\textbf{(1,0 điểm)} Một chiếc ấm điện đun nước siêu tốc vận hành ở điều kiện mạch ổn định với công suất tiêu thụ định mức bằng $P = 1800\text{ W}$. Cho biết các hằng số thông số kĩ thuật đặc trưng của môi chất nước lỏng tinh khiết bao gồm: hằng số nhiệt dung riêng $c = 4,2 \cdot 10^3\text{ J/(kg.K)}$, ẩn nhiệt hóa hơi riêng bão hòa $L = 2,26 \cdot 10^6\text{ J/kg}$ và khối lượng riêng chất lưu $\rho = 1000\text{ kg/m}^3$ (tương đương trị số $1\text{ kg/lít}$).
	\begin{enumerate}
		\item[a,] Tính khoảng thời gian kĩ thuật cần thiết (theo đơn vị giây) để đun nóng $1,5\text{ lít}$ nước lỏng có vạch nhiệt độ ban đầu từ mức $25^\circ\text{C}$ đạt tiến đến khi bắt đầu sôi ở điều kiện áp suất tiêu chuẩn.
		\item[b,] Nếu người sử dụng quên ngắt nguồn, tiếp tục để cho nước lỏng trong ấm duy trì trạng thái sôi bão hòa bùng lên thêm một khoảng thời gian dài đúng bằng 3 phút 26 giây nữa thì lượng nước còn dính lại bên trong ấm điện bấy giờ bằng bao nhiêu lít?
	\end{enumerate}
	\loigiai{
		\begin{enumerate}
			\item[a,] Dựa vào hằng số khối lượng riêng, dung tích nước đem đun $1,5\text{ lít}$ tương ứng khối lượng vật chất $m = 1,5\text{ kg}$. \\
			Ở điều kiện áp suất tiêu chuẩn, nước tinh khiết thực hiện pha sôi ở vạch điểm $t_{\text{sôi}} = 100^\circ\text{C}$. Lượng nhiệt thu vào để nâng khối nước lên điểm sôi đạt mức bằng:
			$$Q_1 = m \cdot c \cdot (t_{\text{sôi}} - t_0) = 1,5 \cdot 4,2 \cdot 10^3 \cdot (100 - 25) = 6300 \cdot 75 = 472500\text{ J}$$
			Do hệ lý tưởng bỏ qua mọi sự thất thoát hao phí, điện năng tiêu thụ chuyển hóa trọn vẹn thành dòng nhiệt đun: $W = Q_1 \Rightarrow P \cdot \tau_1 = Q_1$. \\
			Thời gian đun nước đếm bằng giây cần thiết là:
			$$\tau_1 = \frac{Q_1}{P} = \frac{472500\text{ J}}{1800\text{ W}} = 262,5\text{ s} \quad \text{(Tương đương 4 phút 22,5 giây)}$$
			\item[b,] Quy đổi khoảng thời gian nước lỏng tiếp tục duy trì sôi bùng sang đơn vị giây chuẩn: 
			$$\tau_2 = 3\text{ phút } 26\text{ giây } = 3 \cdot 60 + 26 = 206\text{ s}$$
			Tổng năng lượng nhiệt lượng ấm điện tiếp tục tỏa cấp vào chất lưu ở giai đoạn sôi này bằng:
			$$Q_2 = P \cdot \tau_2 = 1800\text{ W} \cdot 206\text{ s} = 370800\text{ J}$$
			Năng lượng này tiêu hao hoàn toàn phục vụ ẩn nhiệt chuyển pha làm bốc hơi một lượng nước lỏng dời hệ: $Q_2 = \Delta m \cdot L$. \\
			Khối lượng phần hơi nước bão hòa đã bứt phá thoát đi ra ngoài không khí bằng:
			$$\Delta m = \frac{Q_2}{L} = \frac{370800\text{ J}}{2,26 \cdot 10^6\text{ J/kg}} \approx 0,164\text{ kg}$$
			Lượng nước lỏng thực tế còn dư dính lại an toàn bên trong ấm siêu tốc sau quá trình bốc hơi:
			$$m_{\text{còn}} = m - \Delta m = 1,5 - 0,164 = 1,336\text{ kg} \approx 1,34\text{ kg}$$
			Quy đổi tương ứng sang vạch số đo dung tích thể tích: $V_{\text{còn}} \approx 1,34\text{ lít}$.
		\end{enumerate}
	}
\end{ex}

\Closesolutionfile{ans}

\newpage
%%%%%%%%%%%%%%%%%%%%%%%%%%%%%%%%%%%%%%%%%%%%%%%%%%%%%%%%%%%%%%%%%%%%%%%%%
% KHUNG ĐÁP ÁN BẢNG TỰ ĐỘNG IN KẾT XUẤT CỦA HỆ THỐNG
%%%%%%%%%%%%%%%%%%%%%%%%%%%%%%%%%%%%%%%%%%%%%%%%%%%%%%%%%%%%%%%%%%%%%%%%%
\begin{center} \textbf{BẢNG ĐÁP ÁN THAM KHẢO LỚP 12 --- MÃ ĐỀ 429} \end{center}

\noindent \textbf{Phần I (Câu hỏi trắc nghiệm nhiều phương án lựa chọn):}